% Pista ana-25j-q4a · Anàlisi
% Procedència: PAU Catalunya 2025 · Convocatòria de juny · Sèrie 1 · Exercici 4 — Opció A
\documentclass[11pt,a4paper]{article}
\usepackage{pista}
\pistainfo{Catalunya 2025}{Convocatòria de juny}{Sèrie 1}

\begin{document}

\section*{\textcolor{blauPista}{Exercici 4 --- Opció A}}

\noindent La vela major d'un veler té forma semiparabòlica i està delimitada per les gràfiques de $f(x) = -x^{2} + 25$, $y = 0$ i $x = 0$, tal com s'indica a la figura següent:

\noindent\textit{(A la figura: la vela és la regió del primer quadrant compresa entre la paràbola i els eixos de coordenades.)}

\noindent La vela té dues parts separades per la recta $y = 9$. Per a construir-la, s'empra un teixit de niló a la part superior, que costa $50\,\text{\euro}/\text{u}^{2}$, i un teixit de polièster a la part inferior, que costa $70\,\text{\euro}/\text{u}^{2}$. Calculeu el cost total del material que es necessita per a construir aquesta vela. \hspace*{\fill}\textnormal{[2,5~punts]}

\pista{En primer lloc, troba els dos punts de tall que necessites: el de la paràbola amb la recta $y = 9$, resolent $f(x) = 9$, i el de la paràbola amb l'eix d'abscisses, resolent $f(x) = 0$. De cada equació, queda't amb la solució positiva, perquè la vela és al primer quadrant. Anomena-les $x_{1}$ i $x_{2}$, respectivament.

\medskip
A continuació, la part superior és entre la paràbola i la recta $y = 9$, des de l'eix d'ordenades fins a $x_{1}$:
\[
A_{\text{sup}} = \left| \int_{0}^{x_{1}} \bigl(f(x) - 9\bigr)\,dx \right|
\]
Per a la part inferior, et convé calcular primer l'àrea de tota la vela, entre la paràbola i l'eix d'abscisses:
\[
A_{\text{total}} = \left| \int_{0}^{x_{2}} f(x)\,dx \right|
\]
Si a l'àrea de tota la vela li restes $A_{\text{sup}}$, obtindràs l'àrea de la part inferior.

\medskip
Finalment, multiplica cada àrea pel preu del seu teixit i suma els dos resultats.}

\vspace{0.6em}

\end{document}
